\section{Conclusion}\label{sec:conclusion}

This paper develops a comparable, scenario-based test of whether a prospective AI dividend can become usable fiscal space. It carries a conditional technological gain through domestic translation and fiscal capture, then compares the resulting flow with instrument cost, debt consistency, and required persistence on a common denominator. Relative to static redistribution, the contribution is a change in question: not how much a policy redistributes at one point in time, but when and in which country--instrument--scenario cells its projected financing source could sustain a recurrent commitment.

Under the primary channel-based specification, no deterministic crossing occurs in the non-stress scenarios (low, mid, frontier) because domestic translation and fiscal capture are jointly insufficient. Crossings appear only under diagnostic stress, and the debt screen shows why accounting coverage need not constitute usable fiscal space. Measured against the same financing condition, the non-crossings are informative bounds on what the modeled dividend can support.

The central limitation is conditionality: four economies selected for contrast do not form a representative regional panel, and the results depend on the stated scenarios, parameter supports, and uncertainty design rather than identifying causal effects. The model-implied-frequency and derived-horizon qualifications established in Sections~\ref{sec:design} and~\ref{sec:results} therefore bound the comparison as evidence evolves.

The resulting decision rule is restrained. AI-linked fiscal space should not be budgeted as a short-run source for universal social protection. Because the shortfall is joint, governments should strengthen translation and capture together without committing against gains that have not materialized, then apply the debt and persistence screens before recognizing recurrent revenue. Until those conditions are observed, AI-related receipts are upside to be reassessed, while universal commitments remain backed by conventional resources.
